%%%%%%%%%%%%%%%%%%%%%%%%%%%%%%%%%
%
%       EXERCÍCIO
%
%%%%%%%%%%%%%%%%%%%%%%%%%%%%%%%%%

\ifdefstring{\atividade}{prova}{%
    \ifdefstring{\modo}{objetivo}{%
        \renewcommand{\valorquestao}{\ValorQObj\ ponto}
    }{%
        \renewcommand{\valorquestao}{\ValorQDisc\ pontos}
    }%
}%

\begin{exercicioBanco}[\valorquestao]
Julgue as afirmações a seguir, classificando cada uma como verdadeira (V) ou falsa (F).

\begin{enumerate}[\(\bullet\)]
    \item \(\{2,4\} \subseteq \{1,2,3,4,5\}\).

    \item O conjunto \(\{x \mid x\) é consoante da palavra ``banana''\(\}\) é igual ao conjunto \(\{b,n\}\).

    \item \(5 \in \{x \in \bbN \mid x < 5\}\).

    \item \(\{-2\} \subseteq \{x \in \bbN \mid x^2 = 4\}\).

\end{enumerate}

% Define as alternativas
\newcommand{\alternativas}{%
    \begin{center}
        \begin{tabularx}{\textwidth}{XXXXX}
            (a) (V, V, F, F). &
            (b) (V, F, V, F). &
            (c) (F, V, F, V). &
            (d) (V, V, V, F). &
            (e) (F, F, V, V).
        \end{tabularx}
    \end{center}
}

% Define a resposta correta
\newcommand{\resposta}{A}

% Lógica condicional para exibição
\ifdefstring{\atividade}{lista}{%
    \alternativas
    \vspace{0.5em}
    
    \noindent\textbf{Resposta:} letra \textbf{\resposta}.
}{%
    \ifdefstring{\modo}{objetivo}{%
        \alternativas
    }{%
        % modo = discursiva → não mostra alternativas
    }
}
\end{exercicioBanco}

